% !TEX program = xelatex
\documentclass[UTF8,a4paper,12pt,openany]{ctexbook}

\usepackage{amsmath,amssymb,mathtools}
\usepackage{geometry}
\usepackage{booktabs,longtable,array,tabularx,multirow}
\usepackage{graphicx}
\usepackage{xcolor}
\usepackage{listings}
\usepackage{tikz}
\usetikzlibrary{arrows.meta,positioning,shapes.geometric,calc}
\usepackage{enumitem}
\usepackage{fancyhdr}
\usepackage{hyperref}
\usepackage{bookmark}
\usepackage{caption}
\usepackage{float}
\usepackage{siunitx}
\usepackage{etoolbox}

\setmonofont{Consolas}

\geometry{left=2.45cm,right=2.25cm,top=2.45cm,bottom=2.45cm,headheight=15pt}
\setlength{\parindent}{2em}
\setlength{\parskip}{0.25em}
\setlist{nosep,leftmargin=2.2em}
\hypersetup{
  colorlinks=true,
  linkcolor=blue!55!black,
  urlcolor=teal!60!black,
  citecolor=green!40!black,
  pdftitle={SR3机械臂项目二与项目三超详细教程},
  pdfauthor={机器人课程设计小组}
}

\definecolor{codebg}{RGB}{247,248,250}
\definecolor{codekw}{RGB}{0,75,150}
\definecolor{codecomment}{RGB}{80,120,80}
\definecolor{codestring}{RGB}{160,70,40}
\definecolor{boxblue}{RGB}{235,245,255}
\definecolor{boxyellow}{RGB}{255,249,225}
\definecolor{boxred}{RGB}{255,238,238}

\lstdefinelanguage{Matlab}{
  keywords={break,case,catch,classdef,continue,else,elseif,end,for,function,
    global,if,otherwise,parfor,persistent,return,spmd,switch,try,while},
  sensitive=true,
  morecomment=[l]\%,
  morestring=[m]'
}
\lstset{
  language=Matlab,
  basicstyle=\ttfamily\footnotesize,
  keywordstyle=\color{codekw}\bfseries,
  commentstyle=\color{codecomment},
  stringstyle=\color{codestring},
  backgroundcolor=\color{codebg},
  numbers=left,
  numberstyle=\tiny\color{gray},
  numbersep=8pt,
  frame=single,
  rulecolor=\color{gray!35},
  breaklines=true,
  breakatwhitespace=false,
  showstringspaces=false,
  columns=fullflexible,
  keepspaces=true,
  tabsize=4,
  captionpos=b,
  literate={λ}{{$\lambda$}}1 {≥}{{$\geq$}}1
}

\pagestyle{fancy}
\fancyhf{}
\fancyhead[L]{SR3机械臂课程设计}
\fancyhead[R]{项目二：运动学 \quad 项目三：路径规划}
\fancyfoot[C]{\thepage}

\newcommand{\q}{\boldsymbol{q}}
\newcommand{\qd}{\dot{\boldsymbol{q}}}
\newcommand{\qdd}{\ddot{\boldsymbol{q}}}
\newcommand{\R}{\boldsymbol{R}}
\newcommand{\p}{\boldsymbol{p}}
\newcommand{\T}{\boldsymbol{T}}
\newcommand{\J}{\boldsymbol{J}}
\newcommand{\e}{\boldsymbol{e}}
\newcommand{\norm}[1]{\left\lVert #1\right\rVert}
\newcommand{\code}[1]{\texttt{#1}}

\newenvironment{plainbox}[1][知识点]
{\begin{center}\begin{minipage}{0.94\textwidth}
 \colorbox{boxblue}{\parbox{0.965\textwidth}{\textbf{#1}}}\smallskip\par}
{\end{minipage}\end{center}}

\newenvironment{warningbox}[1][特别注意]
{\begin{center}\begin{minipage}{0.94\textwidth}
 \colorbox{boxyellow}{\parbox{0.965\textwidth}{\textbf{#1}}}\smallskip\par}
{\end{minipage}\end{center}}

\newenvironment{dangerbox}[1][容易写错]
{\begin{center}\begin{minipage}{0.94\textwidth}
 \colorbox{boxred}{\parbox{0.965\textwidth}{\textbf{#1}}}\smallskip\par}
{\end{minipage}\end{center}}

\begin{document}

\begin{titlepage}
  \centering
  \vspace*{1.5cm}
  {\Huge\bfseries SR3 六自由度机械臂课程设计\par}
  \vspace{0.8cm}
  {\LARGE\bfseries 项目二：正运动学、逆运动学与雅可比\par}
  \vspace{0.3cm}
  {\LARGE\bfseries 项目三：路径规划、轨迹生成与跟踪控制\par}
  \vspace{1.2cm}
  {\Large 从零开始的超详细原理、推导、代码与测试教程\par}
  \vfill
  \begin{tabular}{rl}
    课程名称：& 机器人课程设计\\[0.35cm]
    机械臂型号：& xMate SR3\\[0.35cm]
    小组成员：& \underline{\hspace{8cm}}\\[0.35cm]
    学号：& \underline{\hspace{8cm}}\\[0.35cm]
    指导教师：& \underline{\hspace{8cm}}\\[0.35cm]
    完成日期：& \today
  \end{tabular}
  \vfill
  {\large 说明：本稿既可作为课程设计报告底稿，也可作为答辩复习讲义。}
\end{titlepage}

\frontmatter
\chapter*{写在最前面：这份文档怎么读}
\addcontentsline{toc}{chapter}{写在最前面：这份文档怎么读}

这份文档故意写得非常慢、非常细。阅读时不要求提前懂矩阵、机器人学或
MATLAB。每遇到一个公式，文档都会尽量回答四个问题：

\begin{enumerate}
  \item 这个量是什么？
  \item 这个公式从哪里来？
  \item 它在机器人运动中表示什么？
  \item 它在代码中为什么要这样写？
\end{enumerate}

本报告严格对应当前仓库中的 SR3 实现。需要提前说明一个事实：

\nolinkurl{inverse_kinematics_analytical.m} 虽然沿用了“解析逆运动学”的文件名，
但当前代码实际采用的是\textbf{多初值阻尼最小二乘数值搜索}，再通过正运动学
反校验和去重得到多组候选解。因此本文会分开讲：

\begin{itemize}
  \item 课程题目要求的 Pieper 分离法和闭式解析逆解思想；
  \item 当前仓库真正运行的 DLS 数值逆解算法；
  \item 两者为什么不同，以及报告和答辩时应如何如实说明。
\end{itemize}

\tableofcontents
\listoffigures
\listoftables

\mainmatter
\chapter{零基础预备知识}

\section{机器人运动学到底在研究什么}

想象桌面上放着一条由六个旋转关节连接起来的机械臂。每个关节都像门轴一样
可以转动。六个关节角写成
\[
\q=[q_1,q_2,q_3,q_4,q_5,q_6].
\]
这里的 $q_i$ 不是坐标点，而是第 $i$ 个电机轴转过的角度。当前项目统一使用
\textbf{弧度}。常见换算为
\[
180^\circ=\pi\ \mathrm{rad},\qquad
1^\circ=\frac{\pi}{180}\ \mathrm{rad}.
\]

机器人学里最常见的两个问题是：

\begin{description}
  \item[正运动学（FK）] 已知六个关节角，计算机械臂末端的位置和朝向。
  \item[逆运动学（IK）] 已知希望末端到达的位置和朝向，反过来求六个关节角。
\end{description}

正运动学像“知道每节手臂怎么弯，算手掌在哪里”；逆运动学像“我想让手掌碰到
杯子，反推肩、肘、腕应该怎么转”。

\section{标量、向量和矩阵}

\subsection{标量}

只有一个数的量称为标量，例如球形障碍物的半径 $r=90\ \mathrm{mm}$。

\subsection{向量}

三个数排成一列可以表示三维空间中的位置：
\[
\p=
\begin{bmatrix}x\\y\\z\end{bmatrix}.
\]
例如 $\p=[480,200,90]^T$ 表示相对于基坐标系，沿 $x$ 方向
\SI{480}{mm}、沿 $y$ 方向 \SI{200}{mm}、沿 $z$ 方向 \SI{90}{mm}。
上标 $T$ 表示转置，即把行向量变成列向量。

向量的长度称为二范数：
\[
\norm{\p}_2=\sqrt{x^2+y^2+z^2}.
\]
代码中的 \code{norm(p)} 就是在算这个量。

\subsection{矩阵}

矩阵可以看成按行列摆放的一组数字。旋转矩阵是一个 $3\times3$ 矩阵：
\[
\R=
\begin{bmatrix}
r_{11}&r_{12}&r_{13}\\
r_{21}&r_{22}&r_{23}\\
r_{31}&r_{32}&r_{33}
\end{bmatrix}.
\]
它的三列分别是新坐标系的 $x$、$y$、$z$ 轴在旧坐标系中指向哪里。

\section{坐标系为什么必不可少}

“末端在 $(100,200,300)$”这句话并不完整，因为必须说明相对于谁测量。
机器人通常给每根连杆建立一个坐标系：
\[
\{0\},\{1\},\ldots,\{6\}.
\]
$\{0\}$ 是固定在机器人底座上的基坐标系，$\{6\}$ 固定在末端。

符号 ${}^0\p_6$ 表示“坐标系 6 的原点，在坐标系 0 中的位置”。
符号 ${}^0\R_6$ 表示“坐标系 6 相对于坐标系 0 的朝向”。

\section{三维旋转}

\subsection{绕三个基本轴旋转}

绕 $x$ 轴旋转 $\alpha$：
\[
\R_x(\alpha)=
\begin{bmatrix}
1&0&0\\
0&\cos\alpha&-\sin\alpha\\
0&\sin\alpha&\cos\alpha
\end{bmatrix}.
\]

绕 $y$ 轴旋转 $\beta$：
\[
\R_y(\beta)=
\begin{bmatrix}
\cos\beta&0&\sin\beta\\
0&1&0\\
-\sin\beta&0&\cos\beta
\end{bmatrix}.
\]

绕 $z$ 轴旋转 $\theta$：
\[
\R_z(\theta)=
\begin{bmatrix}
\cos\theta&-\sin\theta&0\\
\sin\theta&\cos\theta&0\\
0&0&1
\end{bmatrix}.
\]

\begin{warningbox}[矩阵乘法的顺序不能交换]
通常 $\R_x\R_z\ne\R_z\R_x$。先绕 $x$ 转再绕 $z$ 转，与先绕 $z$
再绕 $x$ 转，最终朝向一般不同。因此 MDH 公式中的乘法顺序必须与代码完全一致。
\end{warningbox}

\subsection{旋转矩阵的三个重要性质}

合法旋转矩阵满足
\[
\R^T\R=\boldsymbol{I},\qquad
\R^{-1}=\R^T,\qquad
\det(\R)=1.
\]
第二条非常重要：旋转的逆可以直接用转置得到。项目二中
\[
\R_{36}=\R_{03}^T\R_{06}
\]
正是利用了这一性质。

\section{齐次变换矩阵}

只用 $\R$ 能描述朝向，只用 $\p$ 能描述位置。为了把二者放在一个矩阵里，
机器人学使用 $4\times4$ 齐次变换矩阵：
\[
\T=
\begin{bmatrix}
\R&\p\\
0\ 0\ 0&1
\end{bmatrix}.
\]

若
\[
{}^0\T_1=
\begin{bmatrix}{}^0\R_1&{}^0\p_1\\0&1\end{bmatrix},\quad
{}^1\T_2=
\begin{bmatrix}{}^1\R_2&{}^1\p_2\\0&1\end{bmatrix},
\]
那么从 0 到 2 的变换为
\[
{}^0\T_2={}^0\T_1\,{}^1\T_2.
\]
这就是正运动学循环中不断执行 \code{T = T * A\_i} 的原因。

\section{标准 DH 与改进 DH}

DH 方法用四个参数描述相邻两根连杆：
\[
a,\quad \alpha,\quad d,\quad \theta.
\]
直观解释如下：

\begin{table}[H]
\centering
\begin{tabularx}{0.92\textwidth}{cX}
\toprule
参数 & 含义\\
\midrule
$a$ & 沿某个 $x$ 轴方向的连杆长度\\
$\alpha$ & 相邻两个关节轴之间的扭转角\\
$d$ & 沿关节轴方向的偏移距离\\
$\theta$ & 绕关节轴旋转的角度；转动关节中它含有关节变量\\
\bottomrule
\end{tabularx}
\end{table}

本项目采用改进 DH（MDH）：
\[
{}^{i-1}\T_i=
\R_x(\alpha_{i-1})
\T_x(a_{i-1})
\R_z(\theta_i)
\T_z(d_i).
\]
把四个基本变换相乘，得到代码 \code{mdh\_transform.m} 使用的矩阵：
\[
\begin{bmatrix}
c_\theta&-s_\theta&0&a\\
s_\theta c_\alpha&c_\theta c_\alpha&-s_\alpha&-d s_\alpha\\
s_\theta s_\alpha&c_\theta s_\alpha&c_\alpha&d c_\alpha\\
0&0&0&1
\end{bmatrix}.
\]
其中 $c_\theta=\cos\theta$，$s_\theta=\sin\theta$，其余同理。

\section{关节角偏置 offset}

真实电机零位不一定等于数学模型中 $\theta_i=0$ 的姿态，所以代码使用
\[
\theta_i=q_i+\mathrm{offset}_i.
\]
SR3 的偏置为
\[
\mathrm{offset}=
\left[\pi,\frac{\pi}{2},\frac{\pi}{2},\pi,0,0\right].
\]
例如电机读数 $q_2=0$ 时，MDH 公式实际使用
$\theta_2=\pi/2$。漏掉偏置会让整台机器人模型“从第一步就站错姿势”。

\section{叉乘及其物理意义}

两个三维向量的叉乘为
\[
\boldsymbol{a}\times\boldsymbol{b}=
\begin{bmatrix}
a_yb_z-a_zb_y\\
a_zb_x-a_xb_z\\
a_xb_y-a_yb_x
\end{bmatrix}.
\]
其方向同时垂直于 $\boldsymbol{a}$ 和 $\boldsymbol{b}$。

对转动关节，如果关节轴方向为 $\boldsymbol{z}$，从轴上的一点到末端的向量为
$\p_n-\p_i$，则关节转动产生的瞬时末端线速度方向为
\[
\boldsymbol{z}\times(\p_n-\p_i).
\]
这就是几何雅可比线速度部分的来源。

\section{MATLAB 最低限度语法}

\begin{longtable}{p{0.34\textwidth}p{0.58\textwidth}}
\toprule
代码 & 含义\\
\midrule
\endhead
\code{q(:)'} & 先把数据拉成列向量，再转成行向量，保证形状为 $1\times n$\\
\code{zeros(6,n)} & 创建 $6\times n$ 的全零矩阵\\
\code{eye(4)} & 创建 $4\times4$ 单位矩阵\\
\code{A*B} & 矩阵乘法\\
\code{A.*B} & 对应元素逐个相乘\\
\code{A\textbackslash b} & 解线性方程 $Ax=b$，通常比 \code{inv(A)*b} 稳定\\
\code{T(1:3,4)} & 取齐次矩阵前三行、第四列，即位置向量\\
\code{T(1:3,1:3)} & 取左上角旋转矩阵\\
\code{cross(a,b)} & 三维叉乘\\
\code{norm(x)} & 向量长度；矩阵时需留意所用范数\\
\code{size(A,1)} & 矩阵行数\\
\code{end+1} & 在数组末尾添加新元素\\
\code{struct(...)} & 创建带命名字段的数据结构\\
\bottomrule
\end{longtable}

\begin{plainbox}[读代码的核心方法]
看到一行代码时，先写出每个变量的尺寸。例如
\[
\J:6\times6,\quad \e:6\times1,\quad
\J^T(\J\J^T+\lambda^2I)^{-1}\e:6\times1.
\]
尺寸对得上，代码才可能正确；尺寸对不上，公式一定无法直接运行。
\end{plainbox}

\chapter{项目二：正运动学、逆运动学与雅可比}

\section{项目二要解决的完整问题}

项目二的输入是目标末端位姿
\[
\T_{\mathrm{target}}=
\begin{bmatrix}
\R_{\mathrm{target}}&\p_{\mathrm{target}}\\
0&1
\end{bmatrix},
\]
输出是若干组可能的关节角。它不是一个孤立函数，而是一条验证链：
\[
\q\xrightarrow{\mathrm{FK}}\T
\xrightarrow{\mathrm{IK}}\q'
\xrightarrow{\mathrm{FK}}\T'.
\]
只有当 $\T'\approx\T$ 时，逆解才可信。

\section{SR3 的 MDH 参数}

\begin{table}[H]
\centering
\caption{当前代码采用的 SR3 MDH 参数，长度单位为 mm}
\begin{tabular}{ccccc}
\toprule
$i$ & $a_{i-1}$ & $\alpha_{i-1}$ & $d_i$ & $\theta_i$\\
\midrule
1 & 0   & $0$             & 344   & $q_1+\pi$\\
2 & 0   & $\pi/2$         & 0     & $q_2+\pi/2$\\
3 & 290 & $\pi$           & 0     & $q_3+\pi/2$\\
4 & 0   & $\pi/2$         & 290   & $q_4+\pi$\\
5 & 0   & $-\pi/2$        & 0     & $q_5$\\
6 & 136 & $\pi/2$         & 103.5 & $q_6$\\
\bottomrule
\end{tabular}
\end{table}

\section{正运动学如何一步一步计算}

\subsection{第一步：把关节角变成 MDH 表}

\code{build\_mdh\_table.m} 对每个关节执行
\[
\theta_i=q_i+\mathrm{offset}_i,
\]
然后把 $[a_{i-1},\alpha_{i-1},d_i,\theta_i]$ 放入第 $i$ 行。

\subsection{第二步：每一行变成一个局部变换}

调用
\[
\boldsymbol{A}_i=
\mathrm{mdh\_transform}(a_{i-1},\alpha_{i-1},d_i,\theta_i).
\]
$\boldsymbol{A}_i$ 只描述“从前一坐标系走到当前坐标系”。

\subsection{第三步：把六步串起来}

\[
{}^0\T_6=
{}^0\boldsymbol{A}_1
{}^1\boldsymbol{A}_2
{}^2\boldsymbol{A}_3
{}^3\boldsymbol{A}_4
{}^4\boldsymbol{A}_5
{}^5\boldsymbol{A}_6.
\]
代码从 \code{T = eye(4)} 开始，每次右乘一个 $\boldsymbol{A}_i$。
同时保存中间结果 \code{T\_all(:,:,i+1)}，因为碰撞检测和雅可比不仅需要
末端，还需要每个关节原点的位置。

\section{逆运动学为什么比正运动学难}

正运动学只需按确定顺序代入并相乘，通常只有一个结果。逆运动学困难在于：

\begin{itemize}
  \item 同一末端位姿可能对应多组关节角；
  \item 有些目标超出工作空间，根本无解；
  \item 接近奇异位形时，某些关节运动效果重合；
  \item 三角函数是周期函数，$q$ 与 $q+2k\pi$ 表示相同转角；
  \item 机械臂还受关节限位约束。
\end{itemize}

\section{课程要求的 Pieper 分离法}

\subsection{核心思想}

对后三根关节轴交于一点的六轴球腕机械臂，可以把问题拆开：

\begin{enumerate}
  \item 前三轴负责把腕中心送到正确位置；
  \item 后三轴负责把末端转到正确朝向。
\end{enumerate}

因此先解位置，再解姿态。这称为位置与姿态解耦。

\subsection{腕中心}

在最简单的球腕模型中，末端原点沿末端 $z$ 轴距离腕中心 $d_6$，因此
\[
\p_w=\p_6-d_6\R_6(:,3).
\]
$\R_6(:,3)$ 是末端 $z$ 轴在基坐标系中的方向。

\begin{dangerbox}[当前 SR3 参数并非最简单模板]
当前模型第六行同时含 $a_5=136$ 和 $d_6=103.5$。这意味着末端固定偏置不只
沿 $z_6$ 方向，简单公式 $\p_w=\p-d_6\R(:,3)$ 不足以完整剥离偏置。
这也是当前仓库最终转向通用数值 IK 的重要原因。不能把模板球腕公式不加检查地
宣称为当前代码的精确实现。
\end{dangerbox}

\subsection{前三关节的几何解}

对于理想化平面二连杆，设腕中心为
$[x_w,y_w,z_w]^T$。第一轴先决定腕中心位于机器人哪一侧：
\[
\theta_1=\operatorname{atan2}(y_w,x_w).
\]
另一肩部构型常可写为 $\theta_1+\pi$。

把三维问题投影到关节 2、3 所在平面：
\[
r=x_w\cos\theta_1+y_w\sin\theta_1,\qquad
s=z_w-d_1.
\]
若两段有效长度为 $L_1,L_2$，根据余弦定理：
\[
D=\frac{r^2+s^2-L_1^2-L_2^2}{2L_1L_2}.
\]
当 $|D|>1$ 时三角形无法组成，目标不可达。当 $|D|\le1$ 时：
\[
\theta_3=
\operatorname{atan2}\left(\pm\sqrt{1-D^2},D\right).
\]
正负号产生肘上、肘下两组解。再由
\[
\theta_2=
\operatorname{atan2}(s,r)
-\operatorname{atan2}
\left(L_2\sin\theta_3,L_1+L_2\cos\theta_3\right)
\]
求得肩关节俯仰角。

\subsection{后三关节的姿态解}

前三关节确定后，先算 $\R_{03}$，目标姿态为 $\R_{06}$，于是
\[
\R_{36}=\R_{03}^{-1}\R_{06}=\R_{03}^{T}\R_{06}.
\]
再把 $\R_{36}$ 与腕部三个旋转矩阵的乘积逐元素比较，就能通过
\code{atan2} 提取 $\theta_4,\theta_5,\theta_6$。

一般情况下 $\theta_5$ 的正负分支对应腕部翻转，从而形成两组腕部解。
于是最多有
\[
2\text{（肩）}\times2\text{（肘）}\times2\text{（腕）}=8
\]
组解。

\subsection{腕部奇异}

当 $\sin\theta_5\approx0$ 时，第 4 轴和第 6 轴的效果耦合。此时不是没有解，
而是 $\theta_4$ 与 $\theta_6$ 各自不唯一，通常只有两者之和或差能够确定。
程序必须人为选一个约定，例如令 $\theta_4=0$，再计算 $\theta_6$。

\section{当前代码真正使用的 DLS 数值逆运动学}

\subsection{把逆运动学改写成“不断纠错”}

假设当前猜测为 $\q_k$。用正运动学得到当前位姿 $\T(\q_k)$，与目标比较得到
六维误差：
\[
\e=
\begin{bmatrix}
\e_p\\
\e_o
\end{bmatrix}.
\]
前三维 $\e_p$ 是位置误差：
\[
\e_p=\p_{\mathrm{target}}-\p_{\mathrm{current}}.
\]
后三维 $\e_o$ 是姿态误差。代码先算
\[
\R_{\mathrm{err}}=
\R_{\mathrm{target}}\R_{\mathrm{current}}^T,
\]
再把误差旋转矩阵转换为“旋转轴乘旋转角”的三维向量。

\subsection{小运动近似与雅可比}

在关节角变化很小时：
\[
\Delta\boldsymbol{x}\approx\J(\q)\Delta\q.
\]
希望 $\Delta\boldsymbol{x}$ 抵消当前误差，即
\[
\J\Delta\q\approx\e.
\]
如果直接使用 $\Delta\q=\J^{-1}\e$，会遇到两个问题：$\J$ 可能不可逆，
接近奇异位形时数值会爆炸。

\subsection{阻尼最小二乘}

DLS 使用
\[
\Delta\q=
\J^T\left(\J\J^T+\lambda^2\boldsymbol{I}\right)^{-1}\e.
\]
其中 $\lambda>0$ 是阻尼系数。$\lambda^2I$ 像给病态矩阵加一层缓冲，
使其更容易求解。

代码没有写 \code{inv(A)}，而是：
\begin{lstlisting}
A = J * J' + (opts.lambda^2) * eye(6);
dq = J' * (A \ e);
\end{lstlisting}
因为 \code{A\textbackslash e} 会选择更稳定的线性方程求解方法。

更新规则是
\[
\q_{k+1}=\q_k+\Delta\q^T.
\]
迭代直到 $\norm{\e}$ 小于阈值，或达到最大迭代次数。

\subsection{阻尼系数的影响}

\begin{itemize}
  \item $\lambda$ 太小：远离奇异点时收敛快，但奇异附近可能出现巨大步长；
  \item $\lambda$ 太大：更稳定，但每次只走很小一步，收敛慢且可能残留误差；
  \item 当前单次 IK 默认常用 $0.01$，多初值封装中使用 $0.05$。
\end{itemize}

\section{几何雅可比如何推导}

\subsection{雅可比的尺寸}

六自由度机械臂的几何雅可比为
\[
\J=
\begin{bmatrix}
\J_v\\
\J_\omega
\end{bmatrix}\in\mathbb{R}^{6\times6}.
\]
它把关节速度映射为末端速度：
\[
\begin{bmatrix}\boldsymbol{v}\\\boldsymbol{\omega}\end{bmatrix}
=\J(\q)\dot{\q}.
\]

\subsection{转动关节的一列}

第 $i$ 个转动关节的轴方向为 $\boldsymbol{z}_i$，轴上一点为 $\p_i$，
末端为 $\p_n$。该关节单位角速度引起：
\[
\J_{v,i}=\boldsymbol{z}_i\times(\p_n-\p_i),\qquad
\J_{\omega,i}=\boldsymbol{z}_i.
\]

\subsection{MDH 中为什么不能直接拿累计矩阵的 z 轴}

本项目单步顺序为
\[
\R_x(\alpha)\T_x(a)\R_z(\theta)\T_z(d).
\]
真正的关节转轴是执行前两个固定变换之后、执行 $\R_z(\theta)$ 之前的
$z$ 轴。所以代码先计算
\[
\T_{\mathrm{pre}}=
\T\,\R_x(\alpha_i)\T_x(a_i),
\]
再从 \code{T\_pre(1:3,3)} 取轴方向，从
\code{T\_pre(1:3,4)} 取轴上一点。

\section{姿态误差的轴角表示}

任意旋转矩阵对应一个旋转轴 $\boldsymbol{u}$ 和角度 $\theta$。角度由
\[
\cos\theta=\frac{\operatorname{tr}(\R)-1}{2}
\]
得到。旋转轴乘角度为
\[
\boldsymbol{\omega}=
\frac{\theta}{2\sin\theta}
\begin{bmatrix}
R_{32}-R_{23}\\
R_{13}-R_{31}\\
R_{21}-R_{12}
\end{bmatrix}.
\]
当 $\theta$ 很小时，分母也很小，因此代码直接返回零向量，避免数值除零。

\section{多初值为什么能找到多组解}

DLS 是局部方法：从哪个初始点出发，通常会落入附近的某一个解。
\code{build\_seed\_set} 先准备若干人为设计的肩、肘、腕构型，再在关节限位中
生成 30 组固定随机种子。每个种子独立运行 DLS，最终可能到达不同解。

候选解需经过四道筛选：

\begin{enumerate}
  \item 数值迭代必须声明收敛；
  \item 正运动学回代的位置和姿态误差必须足够小；
  \item 必须满足关节限位；
  \item 与已有解的周期角距离大于阈值，否则视为重复解。
\end{enumerate}

\section{角度归一化与解的选择}

角度具有周期性：
\[
q\equiv q+2k\pi.
\]
代码用
\[
\operatorname{wrap}(x)=\operatorname{mod}(x+\pi,2\pi)-\pi
\]
把角度差折算到 $[-\pi,\pi)$。

选解函数先过滤关节限位，再选择离参考构型最近者：
\[
k^*=\arg\min_k
\norm{\operatorname{wrap}(\q^{(k)}-\q_{\mathrm{ref}})}_2.
\]
这样做的工程意义是减少关节突然大幅跳动，保持动作连续。

\section{当前实现的优点与局限}

\subsection{优点}

\begin{itemize}
  \item 不依赖严格球腕结构，可直接适配当前完整 MDH 模型；
  \item 每个解都经过 FK 反校验，结果可信；
  \item DLS 在奇异附近比普通伪逆稳定；
  \item 多初值可以得到多组候选解。
\end{itemize}

\subsection{局限}

\begin{itemize}
  \item 不能数学保证找到全部解，也不能严格保证最多恰好八组；
  \item 计算量明显高于闭式解析解；
  \item 收敛依赖初值、阻尼和迭代上限；
  \item 位置以 mm、姿态以 rad 混在同一误差向量中，未设置尺度权重；
  \item 接近旋转角 $\pi$ 时，简单轴角公式也可能数值敏感。
\end{itemize}

\begin{warningbox}[报告答辩建议]
不要说“代码已经完整实现 Pieper 闭式解析法”。更准确的表述是：
“报告学习并推导了 Pieper 解耦思想；由于当前 SR3 MDH 末端偏置与简化球腕模板
不完全一致，工程代码采用多初值 DLS，并用 FK 严格反校验候选解。”
\end{warningbox}

\chapter{项目三：路径规划、轨迹生成与跟踪}

\section{路径、轨迹、控制不是同一个东西}

\begin{description}
  \item[路径] 只回答“经过哪些构型”，不包含时间。
  \item[轨迹] 回答“每个时刻处于哪个构型、速度和加速度是多少”。
  \item[控制] 根据期望轨迹和实际状态的误差，产生驱动命令。
\end{description}

项目三的完整流程是
\[
\text{scene}
\xrightarrow{\code{plan\_path}}
\text{path}
\xrightarrow{\code{generate\_trajectory}}
\text{traj}
\xrightarrow{\code{simulate\_tracking}}
\text{result}.
\]

\section{为什么选择关节空间规划}

六轴机械臂的一个状态就是六维向量 $\q$。关节空间规划直接从
$\q_{\mathrm{start}}$ 搜索到 $\q_{\mathrm{goal}}$。

优点是：
\begin{itemize}
  \item 每个搜索节点天然是一组完整关节角；
  \item 容易检查关节限位；
  \item 不需要对每个任务空间采样点反复求 IK；
  \item 多解跳变问题较少。
\end{itemize}

缺点是六维空间难以直观绘制，而且关节空间中的直线映射到任务空间后通常是
弯曲的，仍必须通过正运动学做碰撞检测。

\section{场景数据结构}

\code{scene} 中最重要的字段为：
\[
\q_{\mathrm{start}},\quad
\q_{\mathrm{goal}},\quad
\text{obstacles},\quad T_{\mathrm{total}},\quad dt.
\]
每个球障碍物写成一行
\[
[c_x,c_y,c_z,r].
\]

\section{碰撞检测：为什么不能只看末端}

机械臂末端绕过球，不代表手臂中间的连杆也绕过球。因此代码先通过函数
\nolinkurl{forward_kinematics} 获得所有坐标系原点：
\[
\p_0,\p_1,\ldots,\p_6.
\]
再把相邻点之间的线段近似为机械臂连杆。

\subsection{点到线段最短距离}

线段端点为 $\boldsymbol{A},\boldsymbol{B}$，球心为 $\boldsymbol{C}$。
令
\[
\boldsymbol{AB}=\boldsymbol{B}-\boldsymbol{A}.
\]
把 $\boldsymbol{C}-\boldsymbol{A}$ 投影到 $\boldsymbol{AB}$ 上：
\[
t_0=
\frac{(\boldsymbol{C}-\boldsymbol{A})^T\boldsymbol{AB}}
{\boldsymbol{AB}^T\boldsymbol{AB}}.
\]
如果 $t_0<0$，最近点是 $A$；如果 $t_0>1$，最近点是 $B$；否则最近点在线段
内部。因此
\[
t=\max(0,\min(1,t_0)),
\qquad
\boldsymbol{P}=\boldsymbol{A}+t\boldsymbol{AB}.
\]
最短距离为
\[
d=\norm{\boldsymbol{P}-\boldsymbol{C}}.
\]
若 $d<r$，线段进入球体，判为碰撞。

\subsection{状态合法不等于边合法}

两个端点都不碰撞，中间仍可能穿过障碍。因此对边
$\q_a\rightarrow\q_b$ 做插值：
\[
\q(s)=(1-s)\q_a+s\q_b,\qquad s\in[0,1].
\]
当前规划器取 25 个 $s$ 值逐个检查。采样越密越安全，但计算越慢。

\begin{dangerbox}[离散检测不是连续数学保证]
25 个采样点全部安全，只能说明这些采样点安全。若机械臂在两个采样点之间快速
穿过很小的障碍，仍可能漏检。工程上可减小关节插值步长、按连杆最大位移自适应
采样，或使用连续碰撞检测。
\end{dangerbox}

\section{先试直线为什么合理}

如果起点到终点的关节空间直线本身无碰撞，就没有必要运行随机搜索。
直线路径只需两个路点，路径短、计算快、结果稳定。

因此 \code{plan\_path} 首先调用
\code{is\_edge\_valid(q\_start,q\_goal,...)}。只有直线失败才启动 BiRRT。

\section{RRT 的基本思想}

RRT 是快速扩展随机树。可以把它想成从起点长出一棵树：

\begin{enumerate}
  \item 在关节限位盒子中随机选一个点 $\q_{\mathrm{rand}}$；
  \item 在树中寻找离它最近的节点 $\q_{\mathrm{near}}$；
  \item 从最近节点朝随机点走一个固定步长，得到 $\q_{\mathrm{new}}$；
  \item 如果新状态和连接边都安全，就把新节点加入树；
  \item 重复直到树靠近目标。
\end{enumerate}

\subsection{最近邻}

代码使用普通欧氏距离：
\[
d_k=\norm{\q_k-\q_{\mathrm{rand}}}_2^2.
\]
平方不影响最近节点是谁，同时省去开平方。

\subsection{Steer}

设
\[
\boldsymbol{\delta}=\q_{\mathrm{to}}-\q_{\mathrm{from}}.
\]
若距离小于步长，直接到目标；否则只前进一步：
\[
\q_{\mathrm{new}}=
\q_{\mathrm{from}}+
\Delta q\frac{\boldsymbol{\delta}}{\norm{\boldsymbol{\delta}}}.
\]

\section{为什么使用双向 RRT-Connect}

普通 RRT 只从起点生长。双向方法同时建立：
\[
\mathcal{T}_a\text{ 从起点生长},\qquad
\mathcal{T}_b\text{ 从终点生长}.
\]
一棵树扩出新节点后，另一棵树不断朝它延伸，直到碰撞或成功连接。
每轮再交换两棵树的角色。两边同时长，通常比单向搜索更快相遇。

\begin{figure}[H]
\centering
\resizebox{0.92\textwidth}{!}{%
\begin{tikzpicture}[
  node distance=8mm and 15mm,
  block/.style={rectangle,rounded corners,draw,fill=blue!6,align=center,
    minimum width=34mm,minimum height=8mm},
  decision/.style={diamond,draw,fill=yellow!12,align=center,aspect=2.1},
  line/.style={-{Latex[length=2mm]},thick}
]
\node[block] (start) {读入起点、终点、障碍};
\node[decision,below=of start] (straight) {直线边是否安全？};
\node[block,below left=of straight] (direct) {返回两点路径};
\node[block,below right=of straight] (trees) {建立起点树与终点树};
\node[block,below=of trees] (sample) {随机采样并扩展树 A};
\node[block,below=of sample] (connect) {树 B 尝试连接新节点};
\node[decision,below=of connect] (met) {两棵树连接？};
\node[block,below left=of met] (swap) {交换两棵树继续};
\node[block,below right=of met] (merge) {回溯、拼接、shortcut};
\draw[line] (start)--(straight);
\draw[line] (straight)--node[left]{是}(direct);
\draw[line] (straight)--node[right]{否}(trees);
\draw[line] (trees)--(sample);
\draw[line] (sample)--(connect);
\draw[line] (connect)--(met);
\draw[line] (met)--node[left]{否}(swap);
\draw[line] (swap.west) -| ++(-10mm,0) |- (sample.west);
\draw[line] (met)--node[right]{是}(merge);
\end{tikzpicture}
}
\caption{当前路径规划器的总体流程}
\end{figure}

\section{树的数据结构和路径回溯}

每棵树保存：
\begin{itemize}
  \item \code{tree.q}：每一行是一个节点；
  \item \code{tree.parent}：每个节点的父节点编号。
\end{itemize}

新节点加入时记录它从哪个节点扩展而来。连接成功后，从连接节点沿
\code{parent} 一直回到根，就得到半条路径。两棵树各回溯一次，再翻转并拼接。

\section{目标偏置与确定性}

完全随机采样可能长时间不朝目标前进。当前代码每 5 次迭代直接令
\[
\q_{\mathrm{rand}}=\q_{\mathrm{goal}},
\]
这叫目标偏置。

代码固定 \code{rng(3,'twister')}，所以同一场景通常得到相同结果，便于测试、
调试和报告复现。但固定种子不等于算法对所有新场景都可靠。

\section{当前关键参数为什么这样写}

\begin{table}[H]
\centering
\begin{tabularx}{0.94\textwidth}{lclX}
\toprule
参数 & 当前值 & 单位 & 作用\\
\midrule
最大迭代次数 & 6000 & 次 & 防止搜索无限运行\\
步长 & $0.12\sim0.35$ & rad & 太大易撞，太小搜索慢\\
连接阈值 & 等于步长 & rad & 判断两树是否足够接近\\
边采样数 & 25 & 个 & 检查插值边的碰撞\\
目标偏置 & 每 5 次 & -- & 提高向目标生长概率\\
随机种子 & 3 & -- & 保证结果可复现\\
\bottomrule
\end{tabularx}
\end{table}

\section{Shortcut 路径简化}

RRT 原始路径常有很多锯齿路点。若路径中的第 $i$ 点可以安全直连更后面的第
$j$ 点，就删除二者之间的多余点：
\[
[\q_1,\ldots,\q_i,\q_{i+1},\ldots,\q_j,\ldots]
\rightarrow
[\q_1,\ldots,\q_i,\q_j,\ldots].
\]
当前实现从最远的 $j$ 开始尝试，因此倾向于一次删掉尽可能多的节点。

\section{手工绕行兜底}

若 6000 次随机搜索仍失败，代码围绕中间构型对第 2、3、5 等关节施加
以下几种正负偏移：
\[
\frac{\pi}{6},\quad\frac{\pi}{4},\quad
\frac{\pi}{3},\quad\frac{\pi}{2}.
\]
随后尝试若干三点或四点路径。
这能提高给定课程场景下的成功率，但带有明显启发式，不是通用完备方法。

\section{路径弧长参数化}

得到路点 $\q_1,\ldots,\q_M$ 后，先算每段长度：
\[
\ell_i=\norm{\q_{i+1}-\q_i}_2.
\]
总长度
\[
L=\sum_{i=1}^{M-1}\ell_i.
\]
路点对应的归一化累计弧长为
\[
u_1=0,\qquad
u_i=\frac{\sum_{j=1}^{i-1}\ell_j}{L},\qquad
u_M=1.
\]
这样长路径段会分配更多参数区间，短路径段分配更少。

\section{五次多项式时间缩放从哪里来}

希望标量进度 $s(t)$ 满足六个边界条件：
\[
s(0)=0,\quad s(T)=1,
\]
\[
\dot{s}(0)=\dot{s}(T)=0,
\]
\[
\ddot{s}(0)=\ddot{s}(T)=0.
\]
六个条件至少需要六个系数，所以取五次多项式。令
\[
\tau=\frac{t}{T},\qquad
s(\tau)=a_0+a_1\tau+a_2\tau^2+a_3\tau^3+a_4\tau^4+a_5\tau^5.
\]
代入六个边界条件并解线性方程，得到
\[
a_0=a_1=a_2=0,\quad a_3=10,\quad a_4=-15,\quad a_5=6,
\]
因此
\[
s(\tau)=10\tau^3-15\tau^4+6\tau^5.
\]
导数为
\[
\dot{s}(t)=
\frac{30\tau^2-60\tau^3+30\tau^4}{T},
\]
\[
\ddot{s}(t)=
\frac{60\tau-180\tau^2+120\tau^3}{T^2}.
\]

\section{从标量进度得到每个关节的轨迹}

在某个时刻令 $u=s(t)$，先确定 $u$ 落在哪个路径段
$[u_i,u_{i+1}]$。局部比例为
\[
\lambda=\frac{u-u_i}{u_{i+1}-u_i}.
\]
位置线性插值：
\[
\q(t)=\q_i+\lambda(\q_{i+1}-\q_i).
\]
该段对 $u$ 的导数为
\[
\frac{d\q}{du}=
\frac{\q_{i+1}-\q_i}{u_{i+1}-u_i}.
\]
链式法则给出
\[
\dot{\q}(t)=\frac{d\q}{du}\dot{s}(t),
\qquad
\ddot{\q}(t)=\frac{d\q}{du}\ddot{s}(t).
\]

\begin{warningbox}[当前轨迹并非在路点处完全光滑]
全局 $s(t)$ 很光滑，但 $\q(u)$ 在路点处是分段线性的，方向可能突然改变。
因此 $d\q/du$ 在中间路点处可能跳变，关节速度和加速度也可能不连续。
当前实现满足课程测试的首末边界与数值差分要求，但严格意义上不是全程
$C^2$ 平滑轨迹。改进可使用样条、分段五次多项式或 min-jerk via-point。
\end{warningbox}

\section{PD 控制器}

动态模式控制律为
\[
\boldsymbol{\tau}=
\boldsymbol{K}_p\odot(\q_d-\q)
+\boldsymbol{K}_d\odot(\dot{\q}_d-\dot{\q})
+\ddot{\q}_d.
\]
其中 $\odot$ 表示逐关节相乘。

\begin{itemize}
  \item 比例项像弹簧，位置偏差越大，纠正越强；
  \item 微分项像阻尼，抑制速度误差和振荡；
  \item 前馈加速度提前告诉系统期望怎么加速。
\end{itemize}

当前仿真假设等效惯性矩阵为单位阵，忽略重力、科氏力、摩擦：
\[
\ddot{\q}=\boldsymbol{\tau}.
\]
这只是教学用简化模型，$\tau$ 在这里更像“加速度命令”，不能直接等同真实电机
力矩。

\section{Euler 积分}

每个时间步长 $\Delta t$：
\[
\dot{\q}_{k+1}=\dot{\q}_k+\ddot{\q}_k\Delta t,
\]
\[
\q_{k+1}=\q_k+\dot{\q}_{k+1}\Delta t.
\]
代码先更新速度再更新位置，属于半隐式形式。步长太大会导致数值不稳定。

\section{项目三的局限与改进方向}

\begin{enumerate}
  \item 连杆被近似为无厚度线段，真实机械臂有外壳半径。可把障碍半径扩大安全
  距离，或使用胶囊体碰撞模型。
  \item 随机规划不保证在有限次数内成功。可使用 RRT-Connect 的成熟实现、
  PRM 或 RRT*。
  \item 六个关节直接用同一欧氏度量，未考虑不同关节运动对末端影响不同。
  可设置关节权重。
  \item shortcut 只减少路点，不直接优化安全余量和曲率。
  \item 轨迹未自动检查速度、加速度峰值。可根据限值自动放大 $T$。
  \item 动态模型过于简化。真实控制应考虑
  $M(\q)\qdd+C(\q,\qd)\qd+g(\q)+f(\qd)$。
\end{enumerate}

\chapter{代码阅读、测试判据与答辩准备}

\section{项目二逐文件阅读顺序}

\begin{enumerate}
  \item \code{robot\_params.m}：先确认机器人有几轴、尺寸、偏置和限位。
  \item \code{build\_mdh\_table.m}：理解 $q$ 如何变成 $\theta$。
  \item \code{mdh\_transform.m}：确认采用 MDH，而不是标准 DH。
  \item \code{forward\_kinematics.m}：理解六个局部变换如何累计。
  \item \code{jacobian\_geometric.m}：理解小关节运动如何影响末端。
  \item \code{inverse\_kinematics\_numerical.m}：理解 DLS 单次迭代。
  \item \code{inverse\_kinematics\_analytical.m}：理解多初值搜索、校验和去重。
  \item \code{select\_ik\_solution.m}：理解多解中如何挑一组连续、安全的解。
\end{enumerate}

\section{项目二核心代码逐句解释}

\subsection{正运动学累计}

\begin{lstlisting}
T = eye(4);
for i = 1:params.n
    A_i = mdh_transform(a_i, alpha_i, d_i, theta_i);
    T = T * A_i;
    T_all(:, :, i + 1) = T;
end
\end{lstlisting}

\begin{itemize}
  \item \code{eye(4)} 表示从基坐标系到自身不发生任何变化；
  \item \code{A\_i} 是第 $i$ 节局部变换；
  \item 右乘保持坐标变换链的正确顺序；
  \item 保存 \code{T\_all} 是为了后续使用所有关节位置。
\end{itemize}

\subsection{MDH 关节轴}

\begin{lstlisting}
T_pre = T * rotx4(alpha_i) * transl4(a_i, 0, 0);
z_axis = T_pre(1:3, 3);
p_axis = T_pre(1:3, 4);
Jv = cross(z_axis, p_end - p_axis);
Jw = z_axis;
\end{lstlisting}

\code{T\_pre} 的含义是“已经做完固定的扭转和平移，但还没有执行当前关节的
变量旋转”。因此它的 $z$ 轴就是当前转动关节轴。

\subsection{DLS 更新}

\begin{lstlisting}
ep = T_target(1:3, 4) - T_cur(1:3, 4);
R_err = T_target(1:3, 1:3) * T_cur(1:3, 1:3)';
eo = rotm_to_axisangle(R_err);
e  = [ep; eo];

J = jacobian_geometric(q, params);
A = J * J' + (opts.lambda^2) * eye(6);
dq = J' * (A \ e);
q = q + dq(:)';
\end{lstlisting}

这段代码每轮完成“测误差、算灵敏度、求修正量、更新猜测”四件事。

\section{项目三逐文件阅读顺序}

\begin{enumerate}
  \item \code{scene\_easy.m}、\code{scene\_hard.m}：明确输入是什么。
  \item \code{plan\_path.m} 主函数：先看整体分支。
  \item \code{is\_collision} 和 \code{seg\_point\_dist}：理解安全判据。
  \item \code{extend\_tree}、\code{connect\_tree}：理解 BiRRT。
  \item \code{merge\_trees}、\code{shortcut\_path}：理解路径恢复和简化。
  \item \code{generate\_trajectory.m}：理解时间缩放和链式法则。
  \item \code{pd\_controller.m}、\code{simulate\_tracking.m}：理解理想回放与简化
  动态跟踪的区别。
\end{enumerate}

\section{项目三核心代码逐句解释}

\subsection{直线优先}

\begin{lstlisting}
if is_edge_valid(q_start, q_goal, scene, params, qlim, edge_samples)
    path = [q_start; q_goal];
else
    path = plan_with_birrt(...);
end
\end{lstlisting}

这里不是“偷懒”，而是合理的分层策略：简单问题直接解决，困难问题才使用搜索。

\subsection{扩展一步}

\begin{lstlisting}
idx_near = nearest_node(tree.q, q_target);
q_near = tree.q(idx_near, :);
q_new = steer(q_near, q_target, step);
added = is_state_valid(q_new, ...) && ...
        is_edge_valid(q_near, q_new, ...);
\end{lstlisting}

新节点必须同时满足“点安全”和“从父节点走到它的整条边安全”。

\subsection{轨迹时间缩放}

\begin{lstlisting}
tau = traj.t / scene.T_total;
s = 10*tau.^3 - 15*tau.^4 + 6*tau.^5;
ds_dt = (30*tau.^2 - 60*tau.^3 + 30*tau.^4) ...
        / scene.T_total;
d2s_dt2 = (60*tau - 180*tau.^2 + 120*tau.^3) ...
          / (scene.T_total^2);
\end{lstlisting}

\code{.*} 和 \code{.\^} 表示对时间数组逐元素运算。若误写成矩阵乘方
\code{\^}，MATLAB 会报尺寸错误或得到完全不同的结果。

\section{测试到底在检查什么}

\subsection{IK 自洽测试}

测试不是简单比较 $\q'=\q$，因为逆运动学可以有多解。真正判据是：
\[
\mathrm{FK}(\q')\approx\mathrm{FK}(\q),
\]
并且解集中至少有一组与原始 $\q$ 周期等价。

\subsection{雅可比中心差分测试}

对第 $i$ 个关节分别加减很小扰动 $h$：
\[
\frac{\partial f}{\partial q_i}
\approx\frac{f(\q+h\boldsymbol{e}_i)-f(\q-h\boldsymbol{e}_i)}{2h}.
\]
中心差分误差通常比单边差分小。测试将解析雅可比与数值雅可比的 Frobenius
范数误差比较。

\subsection{路径测试}

依次检查：
\begin{enumerate}
  \item 矩阵尺寸是否为 $M\times6$；
  \item 第一个点是否等于起点；
  \item 最后一个点是否等于终点；
  \item 所有路点是否在关节限位内；
  \item 每条边插值后，所有连杆是否避开所有球。
\end{enumerate}

\subsection{轨迹测试}

检查字段尺寸、等间隔时间、端点位置、端点速度、端点加速度，以及
\[
\frac{\q_{k+1}-\q_k}{dt}
\]
是否与相邻解析速度的平均值基本一致。

\subsection{跟踪测试}

运动学模式直接令实际关节角等于期望关节角，因此理论误差为零。动态模式才会
产生跟踪误差，需要通过 $K_p,K_d$ 调节。

\section{推荐的实验结果图}

课程题目要求的图建议至少包括：
\begin{enumerate}
  \item easy 与 hard 场景三维图：障碍球、起点、终点、中间路点和末端轨迹；
  \item 六个关节角 $q_i(t)$ 曲线；
  \item 六个关节速度 $\dot q_i(t)$ 曲线；
  \item 六个关节加速度 $\ddot q_i(t)$ 曲线；
  \item 机械臂到最近障碍球面的最小安全距离曲线；
  \item 动态模式下关节误差或末端位置误差曲线。
\end{enumerate}

不要凭文字编造数值。应运行当前代码后把真实图和真实测试输出加入最终提交版。

\section{最小安全距离如何计算}

对每个时刻、每个障碍、每根连杆计算线段到球心距离 $d_{j,k}$。球面净空为
\[
c_{j,k}=d_{j,k}-r_k.
\]
该时刻全机械臂最小安全距离：
\[
c_{\min}(t)=\min_{j,k}c_{j,k}(t).
\]
若 $c_{\min}>0$，说明有正安全余量；若等于零，刚好接触；若小于零，发生穿透。

\section{常见报错与定位}

\begin{longtable}{p{0.31\textwidth}p{0.60\textwidth}}
\toprule
现象 & 优先检查\\
\midrule
\endhead
FK 结果完全不对 & MDH/标准 DH 是否混用，offset 是否漏加，矩阵顺序是否写反\\
雅可比位置部分错误 & 是否使用了 MDH 正确关节轴，叉乘顺序是否反了\\
DLS 不收敛 & 初值、阻尼、姿态误差、单位权重、雅可比是否正确\\
IK 有正确位姿但角度不同 & 多解和 $2\pi$ 周期，不应直接逐元素硬比较\\
RRT 一直失败 & 起终点是否已碰撞、步长、边采样、障碍是否封死通道\\
路径端点安全但中间撞球 & 只检查了状态，没有检查边\\
轨迹速度与差分不一致 & 链式法则、段索引、$u$ 节点归一化是否正确\\
动态跟踪振荡 & $K_p$ 过大、$K_d$ 过小、Euler 步长过大\\
\bottomrule
\end{longtable}

\section{答辩可能提问与简洁回答}

\subsection*{为什么用关节空间 RRT，而不是任务空间直线？}
关节空间节点天然满足完整六轴构型表达，容易处理关节限位；任务空间路径需要
频繁求 IK，并可能在不同 IK 分支间跳变。碰撞仍通过 FK 映射到任务空间检查。

\subsection*{为什么用双向 RRT？}
起点树和终点树同时扩展，通常更快连接，尤其适合已知单一起点和终点的问题。

\subsection*{为什么要 shortcut？}
RRT 路径受随机采样影响，节点多且绕。shortcut 用无碰撞长边替换若干短边，
减少路点和路径长度。

\subsection*{为什么五次多项式能保证首末速度、加速度为零？}
因为它的六个系数正好由位置、速度、加速度的六个首末边界条件唯一确定。

\subsection*{DLS 比普通伪逆好在哪里？}
加入 $\lambda^2I$ 后，雅可比接近奇异时求解仍较稳定，不容易产生过大的关节
修正量。

\subsection*{你们的“解析 IK”是真正解析法吗？}
当前运行代码不是严格闭式解析解，而是多初值 DLS 数值法。报告保留 Pieper
分离法作为理论学习，并明确说明当前完整 SR3 MDH 模型使用数值反解和 FK 校验。

\section{AI 工具使用声明模板}

本项目在学习、文档组织和代码解释过程中使用了 OpenAI Codex 辅助工具。
主要用途包括：梳理 MDH、雅可比、DLS、BiRRT 和五次时间缩放的基础原理；
把仓库中的 MATLAB 函数按调用关系整理为可阅读的说明；检查 LaTeX 公式、结构和
编译问题。

具体示例包括：
\begin{enumerate}
  \item 辅助解释为什么 MDH 下雅可比关节轴应在
  $\R_x(\alpha)\T_x(a)$ 之后提取；
  \item 辅助解释 DLS 中阻尼项的数值作用，以及为什么使用 MATLAB
  反斜杠而不是显式求逆；
  \item 辅助整理 BiRRT、碰撞检测、shortcut 和五次时间缩放的逐步说明。
\end{enumerate}

最终提交前，小组成员应自行核对公式与当前代码的一致性，亲自运行测试、生成
实验图、记录参数与结果，并能够在现场不依赖 AI 独立解释算法和代码。算法选择、
实际调参、测试结果和答辩内容应由小组成员共同确认。


\appendix
\chapter{核心 MATLAB 源码附录}

本附录通过 \code{\textbackslash lstinputlisting} 直接读取仓库中的当前源码。
因此代码更新后重新编译，附录也会同步更新。

\section{项目二核心代码}

\subsection{机器人参数}
\lstinputlisting[caption={robot\_params.m}]{../project2/robot_params.m}

\subsection{MDH 参数表构造}
\lstinputlisting[caption={build\_mdh\_table.m}]{../project2/build_mdh_table.m}

\subsection{单节 MDH 变换}
\lstinputlisting[caption={mdh\_transform.m}]{../project2/mdh_transform.m}

\subsection{正运动学}
\lstinputlisting[caption={forward\_kinematics.m}]{../project2/forward_kinematics.m}

\subsection{几何雅可比}
\lstinputlisting[caption={jacobian\_geometric.m}]{../project2/jacobian_geometric.m}

\subsection{DLS 数值逆运动学}
\lstinputlisting[caption={inverse\_kinematics\_numerical.m}]{../project2/inverse_kinematics_numerical.m}

\subsection{多初值逆运动学封装}
\lstinputlisting[caption={inverse\_kinematics\_analytical.m}]{../project2/inverse_kinematics_analytical.m}

\subsection{逆解选择}
\lstinputlisting[caption={select\_ik\_solution.m}]{../project2/select_ik_solution.m}

\section{项目三核心代码}

\subsection{路径规划}
\lstinputlisting[caption={plan\_path.m}]{../project3/plan_path.m}

\subsection{轨迹生成}
\lstinputlisting[caption={generate\_trajectory.m}]{../project3/generate_trajectory.m}

\subsection{PD 控制器}
\lstinputlisting[caption={pd\_controller.m}]{../project3/pd_controller.m}

\subsection{跟踪仿真}
\lstinputlisting[caption={simulate\_tracking.m}]{../project3/simulate_tracking.m}

\chapter{运行命令与提交前检查表}

\section{MATLAB 运行方式}

\begin{lstlisting}
addpath('../project1');
addpath('../project2');
addpath('../project3');

cd('../project2');
main_project2_test;

cd('../project3');
main_project3_test;
\end{lstlisting}

\section{LaTeX 编译方式}

在 \code{latex\_detailed\_report} 文件夹中使用 XeLaTeX：
\begin{lstlisting}[language={}]
xelatex main.tex
xelatex main.tex
\end{lstlisting}
编译两次是为了让目录、页码和交叉引用稳定。

\section{提交前检查表}

\begin{enumerate}
  \item 封面小组成员、学号、教师和日期已填写；
  \item 文档中“数值 IK 与解析 IK 的区别”表述真实；
  \item easy、hard 场景均已实际运行；
  \item 测试输出来自本机真实运行，没有手工伪造；
  \item 3D 场景图和 $q,\dot q,\ddot q$ 曲线已加入；
  \item 最小安全距离曲线已加入；
  \item 速度、加速度峰值与允许值做了比较；
  \item AI 工具使用声明按实际情况修改；
  \item 每位成员能解释 DLS、碰撞检测、BiRRT 和五次时间缩放；
  \item 最终 PDF 中目录、公式、中文、代码换行均正常。
\end{enumerate}


\backmatter
\chapter*{结语}
\addcontentsline{toc}{chapter}{结语}
项目二解决“给定关节角，末端在哪里”和“给定末端位姿，关节应该转多少”；
项目三解决“从一个构型到另一个构型，怎样避开障碍并平稳地走过去”。
把这两部分连起来，就得到一条完整的机器人运动链：
\[
\text{建模}\rightarrow\text{正/逆运动学}\rightarrow
\text{碰撞检测}\rightarrow\text{路径规划}\rightarrow
\text{轨迹生成}\rightarrow\text{闭环跟踪}.
\]
真正需要掌握的不是背下某一行 MATLAB，而是知道每一行代码对应上面链条中的
哪一个物理问题、数学问题和工程约束。

\end{document}
